% Part: normal-modal-logic
% Chapter: axioms-systems
% Section: provability-properties

\documentclass[../../../include/open-logic-section]{subfiles}

\begin{document}

\olfileid{nml}{prf}{prp}

\olsection{خصائص \usetoken{S}{derivability}}

\begin{prop}\ollabel{prop:derivabilityfacts}
  إذا كان~$\Sigma$ نسقًا جهويًا وكانت~$\Gamma$ مجموعة من ال!!{formula}s
  الجهوية، فإن قابلية الاشتقاق تحقق الخصائص الآتية:
  \begin{enumerate}
  \item\ollabel{prop:derivabilityfacts-monotonicity} \emph{الرتابة}: إذا
    كان~$\Gamma \Proves[\Sigma] !A$ و$\Gamma \subseteq \Delta$، فإن
    $\Delta \Proves[\Sigma] !A$؛
  \item\ollabel{prop:derivabilityfacts-reflexivity}
    \emph{الانعكاسية}: إذا كان~$!A \in \Gamma$، فإن
    $\Gamma \Proves[\Sigma] !A$؛
  \item\ollabel{prop:derivabilityfacts-cut} \emph{القطع}: إذا كان
    $\Gamma \Proves[\Sigma] !A$ و$\Delta \cup \{!A\}
    \Proves[\Sigma] !B$، فإن~$\Gamma \cup \Delta \Proves[\Sigma] !B$؛
  \item\ollabel{prop:derivabilityfacts-deduction} \emph{مبرهنة الاستنتاج}:
    $\Gamma \cup \{!B\} \Proves[\Sigma] !A$ إذا وفقط إذا كان
    $\Gamma \Proves[\Sigma] !B \lif !A$؛
  \item \ollabel{prop:derivabilityfacts-ruleT}% \emph{Rule T}: If
    إذا كان~$\Gamma \Proves[\Sigma] !A_1$ و\dots{} و
    $\Gamma \Proves[\Sigma] !A_n$، وكانت
    $!A_1 \lif (!A_2 \lif \cdots (!A_n \lif !B)\cdots)$ مثال إحلال
    لقضية صادقة صدقًا تحليليًا، فإن~$\Gamma \Proves[\Sigma] !B$.
  \end{enumerate}
\end{prop}

وإثبات هذه الخصائص تمرين يسير. ولبيان الانتفاع بالجزء
\olref{prop:derivabilityfacts-ruleT} من~\olref{prop:derivabilityfacts}،
نأخذ المثال الآتي: إذا ثبت~$\Gamma \Proves[\Sigma] !A \lor !B$ وثبت
$\Gamma \Proves[\Sigma] \lnot !A$، لزم~$\Gamma \Proves[\Sigma] !B$.
واختصارًا للعبارة سنكتب~$\Gamma, !A \Proves[\Sigma] !B$ عوضًا عن
$\Gamma \cup \{!A\} \Proves[\Sigma] !B$.

\begin{defn}
  يقال للمجموعة~$\Gamma$ إنها \emph{مغلقة استنتاجيًا} بالنسبة إلى النسق
  $\Sigma$ إذا كانت مشتملة على كل ما يشتق منها؛ أي إذا وفقط إذا كان
  $\Gamma \Proves[\Sigma] !A$ يقتضي $!A \in \Gamma$.
\end{defn}

\end{document}
